% 构造表：4行3列；两个跨格对象各合并2格，得到10个逻辑单元。
\begin{tikzpicture}[x=1cm,y=1cm,font=\figurefont\small,
  arr/.style={-{Stealth[length=2mm]},BookInk,line width=0.9pt}]
  \node[anchor=west,font=\figurefont\bfseries\small] at (0,4.0)
    {行列交叉与跨格对象};
  \node[anchor=west,font=\figurefont\bfseries\small] at (9,4.0)
    {恢复后的逻辑表};
  \foreach \dx in {0,9}{
    \begin{scope}[xshift=\dx cm]
      \fill[BookBlue!12] (0,1.6) rectangle (2,3.2);
      \fill[BookTeal!12] (2,2.4) rectangle (6,3.2);
      \draw[BookInk,line width=0.9pt] (0,0) rectangle (6,3.2);
      \draw[BookRule,line width=0.5pt] (0,0.8) -- (6,0.8);
      \draw[BookRule,line width=0.5pt] (0,1.6) -- (6,1.6);
      \draw[BookRule,line width=0.5pt] (2,0) -- (2,3.2);
      \node[fill=BookBlue!12,inner sep=2pt] at (1,2.4) {项目};
      \node[fill=BookTeal!12,inner sep=2pt] at (4,2.8) {金额（元）};
      \foreach \x/\y/\v in {
        3/2/1月,5/2/2月,1/1.2/甲,3/1.2/12,5/1.2/18,
        1/.4/乙,3/.4/8,5/.4/10}
        \node at (\x,\y) {\v};
    \end{scope}
  }
  \draw[BookMuted,dashed,line width=0.5pt] (4,0) -- (4,3.2);
  \draw[BookMuted,dashed,line width=0.5pt] (0,2.4) -- (6,2.4);
  \draw[BookBlue,line width=1.1pt] (0,1.6) rectangle (2,3.2);
  \draw[BookTeal,line width=1.1pt] (2,2.4) rectangle (6,3.2);
  \node[fill=BookBlue!12,inner sep=2pt] at (1,2.4) {项目};
  \node[fill=BookTeal!12,inner sep=2pt] at (4,2.8) {金额（元）};
  \draw[BookRule,line width=0.5pt] (13,0) -- (13,2.4);
  \draw[BookRule,line width=0.5pt] (11,2.4) -- (15,2.4);
  \draw[arr] (6.3,1.6) -- (8.6,1.6);
  \node[align=center] at (7.45,2.35) {按跨格对象\\合并};
  \node at (3,-0.6) {12个基本格};
  \node at (12,-0.6) {10个逻辑单元};
\end{tikzpicture}
